\documentclass[10pt,a4paper]{amsart}
\usepackage{amsmath,amssymb,amsthm,graphicx}
\usepackage{verbatim}
\newenvironment{reference}{\comment}{\endcomment}
%\newenvironment{reference}{}{}
\newenvironment{slogan}{\comment}{\endcomment}
\newenvironment{history}{\comment}{\endcomment}

% For commutative diagrams we use Xy-pic
\usepackage[all]{xy}

% We use 2cell for 2-commutative diagrams.
\xyoption{2cell}
% Theorem environments.
%
\theoremstyle{plain}
\newtheorem{theorem}[subsection]{Theorem}
\newtheorem{proposition}[subsection]{Proposition}
\newtheorem{lemma}[subsection]{Lemma}

\theoremstyle{definition}
\newtheorem{definition}[subsection]{Definition}
\newtheorem{example}[subsection]{Example}
\newtheorem{exercise}[subsection]{Exercise}
\newtheorem{situation}[subsection]{Situation}

\theoremstyle{remark}
\newtheorem{remark}[subsection]{Remark}
\newtheorem{remarks}[subsection]{Remarks}

\numberwithin{equation}{subsection}

% Macros
%
\def\lim{\mathop{\mathrm{lim}}\nolimits}
\def\colim{\mathop{\mathrm{colim}}\nolimits}
\def\Spec{\mathop{\mathrm{Spec}}}
\def\Hom{\mathop{\mathrm{Hom}}\nolimits}
\def\Ext{\mathop{\mathrm{Ext}}\nolimits}
\def\SheafHom{\mathop{\mathcal{H}\!\mathit{om}}\nolimits}
\def\SheafExt{\mathop{\mathcal{E}\!\mathit{xt}}\nolimits}
\def\Sch{\mathit{Sch}}
\def\Mor{\mathop{\mathrm{Mor}}\nolimits}
\def\Ob{\mathop{\mathrm{Ob}}\nolimits}
\def\Sh{\mathop{\mathit{Sh}}\nolimits}
\def\NL{\mathop{N\!L}\nolimits}
\def\CH{\mathop{\mathrm{CH}}\nolimits}
\def\proetale{{pro\text{-}\acute{e}tale}}
\def\etale{{\acute{e}tale}}
\def\QCoh{\mathit{QCoh}}
\def\Ker{\mathop{\mathrm{Ker}}}
\def\Im{\mathop{\mathrm{Im}}}
\def\Coker{\mathop{\mathrm{Coker}}}
\def\Coim{\mathop{\mathrm{Coim}}}

% Boxtimes
%
\DeclareMathSymbol{\boxtimes}{\mathbin}{AMSa}{"02}

%
% Macros for moduli stacks/spaces
%
\def\QCohstack{\mathcal{QC}\!\mathit{oh}}
\def\Cohstack{\mathcal{C}\!\mathit{oh}}
\def\Spacesstack{\mathcal{S}\!\mathit{paces}}
\def\Quotfunctor{\mathrm{Quot}}
\def\Hilbfunctor{\mathrm{Hilb}}
\def\Curvesstack{\mathcal{C}\!\mathit{urves}}
\def\Polarizedstack{\mathcal{P}\!\mathit{olarized}}
\def\Complexesstack{\mathcal{C}\!\mathit{omplexes}}
% \Pic is the operator that assigns to X its picard group, usage \Pic(X)
% \Picardstack_{X/B} denotes the Picard stack of X over B
% \Picardfunctor_{X/B} denotes the Picard functor of X over B
\def\Pic{\mathop{\mathrm{Pic}}\nolimits}
\def\Picardstack{\mathcal{P}\!\mathit{ic}}
\def\Picardfunctor{\mathrm{Pic}}
\def\Deformationcategory{\mathcal{D}\!\mathit{ef}}
\usepackage{cleveref}
\crefname{originalreference}{original reference}{original references}
\Crefname{originalreference}{Original reference}{Original references}
\makeatletter
\newlabel{example-dgm-dg-cat}{{09LC}{source}{}{sourcepointer-1}{}}
\newlabel{example-dgm-dg-cat@cref}{{[originalreference][0][]09LC}{[1][0][]source}{}{}{}}
\makeatother
\crefname{originalreference}{original reference}{original references}
\Crefname{originalreference}{Original reference}{Original references}
\makeatletter
\newlabel{lemma-derived-restriction}{{09LI}{source}{}{sourcepointer-2}{}}
\newlabel{lemma-derived-restriction@cref}{{[originalreference][0][]09LI}{[1][0][]source}{}{}{}}
\makeatother
\crefname{originalreference}{original reference}{original references}
\Crefname{originalreference}{Original reference}{Original references}
\makeatletter
\newlabel{lemma-hom-derived}{{09KY}{source}{}{sourcepointer-3}{}}
\newlabel{lemma-hom-derived@cref}{{[originalreference][0][]09KY}{[1][0][]source}{}{}{}}
\makeatother
\title{Derived Hom from a K-projective differential graded module is dg-Hom}
\author{The Stacks Project Authors}
\date{stacks-project:dga.tex:local-snapshot-20261012; GNU FDL 1.2 or later}
\begin{document}
\maketitle
\noindent\textit{Editorial note.} All delivered mathematical text is the author's own, copied from the pinned Stacks Project chapter file sources/stacks-dga/pinned-original-source.tex; no statement or proof step was rewritten, reordered or removed.\par\medskip
\noindent\textit{Editorial note.} The counted proof is the author's complete argument: dependencies carried as navigation stubs: Lemma lemma-derived-restriction (construction), Example example-dgm-dg-cat (degree-wise Hom\_K identification), Lemma lemma-hom-derived (property (P) case); all in-source.\par\medskip
\noindent\textit{Editorial note.} The label lemma-compute-derived-restriction names the counted statement, and the proof environment immediately following it in the source is delivered verbatim as its proof.\par\medskip
\noindent\textit{Editorial note.} The counted statement is carried by the author's own lemma environment in the chapter source.\par\medskip
\noindent\textit{Allowed standard prerequisites (editorial).}\begin{itemize}
\item Construction of derived restriction. Same unit, Lemma lemma-derived-restriction. By construction, RHom(N, M) is computed by choosing a quasi-isomorphism M -> I with I a differential graded B-module of property (I) and setting RHom(N, M) = Hom in Mod\textasciicircum{}dg\_(B,d) from N to I.
\item Hom-complex cohomology identification. Same unit, Example example-dgm-dg-cat and Lemma lemma-hom-derived. The cohomology in each degree of the Hom-complex between differential graded modules is Hom in the homotopy category after shifting (standard dg-category fact recorded in the example); the property (P) case of the hypothesis is supplied by Lemma lemma-hom-derived.
\end{itemize}\medskip
\ifdefined\hypertarget\hypertarget{sourcepointer-1}{}\fi
\noindent\textit{Source reference.} Original reference 09LC: sources/stacks-dga/pinned-original-source.tex, source lines 3525--3525. Prerequisite pointer invoked at proof line(s) 5407; it names a standard fact of the same chapter that the counted argument uses, and every step of the counted proof built on it is retained verbatim in the delivered body.\par\medskip
\ifdefined\hypertarget\hypertarget{sourcepointer-2}{}\fi
\noindent\textit{Source reference.} Original reference 09LI: sources/stacks-dga/pinned-original-source.tex, source lines 5272--5272. Prerequisite pointer invoked at proof line(s) 5396; it names a standard fact of the same chapter that the counted argument uses, and every step of the counted proof built on it is retained verbatim in the delivered body.\par\medskip
\ifdefined\hypertarget\hypertarget{sourcepointer-3}{}\fi
\noindent\textit{Source reference.} Original reference 09KY: sources/stacks-dga/pinned-original-source.tex, source lines 2850--2850. Prerequisite pointer invoked at proof line(s) 5409; it names a standard fact of the same chapter that the counted argument uses, and every step of the counted proof built on it is retained verbatim in the delivered body.\par\medskip
\begin{lemma}
\label{lemma-compute-derived-restriction}
Let $R$ be a ring. Let $(A, \text{d})$ and $(B, \text{d})$ be
differential graded $R$-algebras. Let $N$ be a differential
graded $(A, B)$-bimodule. If
$\Hom_{D(B, \text{d})}(N, N') = \Hom_{K(\text{Mod}_{(B, \text{d})})}(N, N')$
for all $N' \in K(B, \text{d})$, for example if $N$
has property (P) as a differential graded $B$-module, then
$$
R\Hom(N, M) = \Hom_{\text{Mod}^{dg}_{(B, \text{d})}}(N, M)
$$
functorially in $M$ in $D(B, \text{d})$.
\end{lemma}

\begin{proof}
By construction (Lemma \ref{lemma-derived-restriction})
to find $R\Hom(N, M)$ we choose a quasi-isomorphism
$M \to I$ where $I$ is a differential graded $B$-module
with property (I) and we set
$R\Hom(N, M) = \Hom_{\text{Mod}^{dg}_{(B, \text{d})}}(N, I)$.
By assumption the map
$$
\Hom_{\text{Mod}^{dg}_{(B, \text{d})}}(N, M) \longrightarrow
\Hom_{\text{Mod}^{dg}_{(B, \text{d})}}(N, I)
$$
induced by $M \to I$ is a quasi-isomorphism, see discussion in
Example \ref{example-dgm-dg-cat}. This proves the lemma.
If $N$ has property (P) as a $B$-module, then we see that the
assumption is satisfied by Lemma \ref{lemma-hom-derived}.
\end{proof}

\end{document}
